\section{A finite interpreter in the two-sided tape model}
\label{sv:interpreter}
Before serialization, delete all outgoing rows from designated halt
states.  As in the recorder constructions, complete every missing
state--symbol pair at a nonhalting state by an unchanged-symbol,
zero-displacement move to a fresh halt state with no outgoing rows.
This finite normalization preserves the halting question, including an
initially halted machine.  In the resulting table, failure to find a
matching row is equivalent to being in a halt state.

Give the interpreted states and symbols positive integer indices, with
blank index one.  Write an index $m$ as a block of $m$ ones followed
by a separator.  Use distinct syntax symbols for begin, row-end,
table-start, table-end, and the three head motions.  Serialize a machine
and input by its initial state, its input-symbol fields in reverse word
order, and its finite list of transition rows.  A row contains current
state, scanned symbol, next state, written symbol, and motion.
Correctness is required on these valid finite descriptions.

First construct a fixed five-tape interpreter.  The tapes hold the
unchanged description, a current-state record, a left-cell stack, a
right-cell stack, and a scratch record.  Each work tape has a sentinel
at cell zero.  A record is a finite unary block immediately to its right,
with its head returned to the sentinel between uses.  A stack consists
of finitely many unary symbol blocks to the right of its sentinel,
written bottom to top; each block-end symbol also carries an origin
bit.  The head rests at its top block-end, or at its sentinel when
empty.  Below the stored blocks the logical stack consists of unmarked
blanks.

The needed routines have finite control and terminating scans.  To copy a
field, erase the old record, return to its sentinel, and copy the field's
ones while scanning it.  To compare a field and a record, scan their
ones in parallel, remember whether they end together, and return the
record head.  To push, copy the scratch unary block after the top and
append the chosen origin-bit block-end.  To pop, retain the top origin
bit in control, erase its end, and traverse its ones backward, copying
one to scratch for every erased one until the previous end or sentinel.
An empty stack instead supplies blank index one and origin bit zero.
All these scans stop on explicitly present finite delimiters.

Initialize the state record from the description.  Push the reversed
input fields onto the right stack, then mark its top block as origin;
for an empty input push a marked blank.  Leave the left stack empty.
In each iteration pop the right top into scratch, remembering its bit.
Search the table for a row matching the state and scratch records.
On a match update the state and scratch to the next state and written
symbol.  A right move pushes the written symbol onto the left; a stay
pushes it back onto the right.  A left move first pushes the written
symbol right, then pops the left (supplying an implicit blank if empty)
and pushes that symbol right.  Every push retains the appropriate
origin bit.  Return the description head to table-start and repeat.
If no row matches, restore the popped symbol and halt; by the normalization
above, this also handles a halt on the first iteration.  The two logical
stacks are therefore exactly the interpreted tape halves, and the
origin bit recovers its absolute head position at every iteration.

These routines compile to a finite five-tape table: instruction location,
comparison outcome, origin bit, and chosen direction have finitely many
values; unbounded interpreted indices stay in unary on tape.  Each
character test, write, copy, or move is an ordinary finite transition,
and the finitely many call sites can be inlined.  Malformed descriptions
may receive arbitrary defaults.

Finally encode the five aligned tracks on one two-sided tape between
left and right fences, with one head bit per track in the product
alphabet.  At a simulated step, sweep the finite active interval to
collect all five scanned symbols into finite control.  Choose the
transition before changing any track.  Then sweep to each head bit in
turn, make that track's write and move, preserving every other track.
If a head crosses a fence, replace that fence by a blank tuple with the
new head bit and put the new fence one cell farther out.  Return to the
left fence between sweeps.  Only finitely many tracks, symbols, and
pending transition choices are involved, so these are again finite
rules.  Every sweep terminates, and all data and head bits match the
five-tape machine at step boundaries.  Halt exactly when it halts.
This yields a fixed universal single-tape table.  Serialization and
fluid compilation use only the finite description and input, and never
run the interpreted computation in advance.

The same fence construction also converts any fixed finite number of
tapes to one before applying a recorder: collect the finitely many head
symbols, move the fences outward when needed, update one track at a time,
and return to the left fence.  If a one-sided tape convention is desired,
add a preserved bit at its left endpoint and implement that convention's
left-end behavior in the marked-symbol transitions.  Both reductions
preserve the initialized halting question without changing the geometric
arguments.
